\subsection{Sprint 2}

A Sprint 2 começou em um cenário que ainda carregava problemas da entrega
anterior. Embora o time tivesse avançado no desenvolvimento, permaneciam
\textit{pull requests} atrasados e tarefas que dependiam de trabalhos ainda não
integrados. Além disso, esta sprint teve menos tempo disponível para a entrega e
coincidiu com o período de provas de grande parte dos integrantes, reduzindo a
capacidade de trabalho do grupo justamente quando ainda havia pendências
acumuladas.

No dia 22 de setembro, um dia após o início da sprint, concentrei minha atuação
na correção de PRs do \textit{backend} e do \textit{frontend}. Esse trabalho deu
continuidade ao esforço que eu já havia realizado no final da Sprint 1, quando os
AGES IV precisaram revisar e integrar uma fila de contribuições pendentes. As
correções não se limitavam a aprovar os PRs: era necessário verificar o código,
identificar problemas que impediam a integração e devolver ou ajustar as
contribuições para que pudessem compor uma versão minimamente estável da
aplicação.

A combinação do prazo menor, da semana de provas e do acúmulo de tarefas
comprometeu novamente o andamento da sprint. Parte do trabalho permaneceu
bloqueada por dependências entre tarefas, alguns PRs demoraram a ser corrigidos e
as informações sobre o progresso nem sempre estavam atualizadas. Com isso, uma
quantidade significativa de trabalho ficou concentrada no final do período. No
dia anterior à apresentação, durante o \textit{code freeze}, nós, alunos AGES IV,
precisamos revisar, corrigir e integrar às pressas o que ainda poderia ser
aproveitado para a entrega.

Mesmo com esse esforço, a entrega ficou incompleta e com qualidade inferior à
que pretendíamos. Nem todas as funcionalidades planejadas estavam disponíveis,
e a versão em produção não pôde ser apresentada. A retrospectiva realizada no
dia da entrega registrou, entre os principais problemas percebidos pelo time, as
tarefas bloqueadas, a demora na revisão dos PRs, a concentração do trabalho no
final da sprint, a falta de comunicação entre as squads e a ausência de
informações suficientes nas tarefas. Ao mesmo tempo, o grupo reconheceu que os
integrantes procuraram se manter ativos apesar das provas e que os AGES IV
conseguiram finalizar a apresentação em meio à correria.

Um problema que passou a aparecer com mais força nessa retrospectiva foi a
organização do quadro do ClickUp. O quadro havia sido configurado de modo que os
cartões principais representavam as histórias de usuário, enquanto as tarefas de
desenvolvimento ficavam registradas como subtarefas dentro delas. Na prática,
essa estrutura dificultava visualizar quem estava responsável por cada atividade,
quais itens estavam bloqueados e o que precisava ser priorizado. Também tornava
mais trabalhoso mover e atualizar os cartões durante o acompanhamento da sprint.
Uma organização com uma visão para as histórias de usuário e outra para as
tarefas individuais teria permitido acompanhar os dois níveis de trabalho sem
esconder as atividades executadas dentro das histórias.

\medskip

\noindent\fbox{%
    \parbox{0.95\linewidth}{%
        \setlength{\parindent}{15pt}
        \itshape A Sprint 2 mostrou que o esforço realizado
        no final da Sprint 1 resolveu a entrega imediata, mas não eliminou as
        causas do acúmulo de trabalho. Voltamos a chegar ao \textit{code freeze}
        com PRs atrasados, tarefas bloqueadas e pouca clareza sobre o andamento
        real da sprint. O período de provas agravou o problema, mas não explica
        sozinho a baixa qualidade da entrega.

        Como AGES IV, reconheço que nossa atuação foi novamente mais reativa do
        que preventiva. Conseguimos reunir o grupo e fazer as correções urgentes
        no dia anterior à apresentação, mas esse esforço de última hora não
        substituiu o acompanhamento que deveria ter ocorrido ao longo de toda a
        sprint. Eu participei da correção dos PRs, porém poderia ter cobrado antes
        a atualização das tarefas, identificado os bloqueios e ajudado a
        redistribuir o trabalho antes que ele se acumulasse.

        O ClickUp também deixou de ser apenas uma ferramenta pouco atualizada e
        passou a representar um problema de organização. A estrutura escolhida
        escondia as tarefas individuais dentro das histórias de usuário e
        dificultava enxergar o fluxo de trabalho. Tentei melhorar esse quadro ao
        longo do semestre, mas não consegui estabelecer uma organização que o
        time realmente adotasse. Essa experiência reforçou para mim que uma
        ferramenta de gestão não resolve a falta de acompanhamento por si só: sua
        estrutura precisa representar o modo como o time trabalha, e todos
        precisam manter as informações atualizadas.
    }%
}
